%% ma: L12-B14-0129
%% lop: 12
%% chuong: 5
%% bai: 14
%% dang: 12.14.25
%% loai: TN4
%% muc: TH
%% dap_an: "A"
%% nguon: "Ảnh giáo viên gửi 25/11/2026 (tài liệu ôn tập Phương trình mặt phẳng), Phần I Câu 12"
%% kien_thuc: "Bài 7 (tọa độ vectơ), Bài 14 (phương trình mặt phẳng đi qua một điểm và có vectơ pháp tuyến)"
%% trang_thai: chinh-thuc
\begin{ex}
Trong không gian $Oxyz$, cho hai điểm $A(0;1;1)$ và $B(1;2;3)$. Viết phương trình của mặt phẳng $(P)$ đi qua $A$ và vuông góc với đường thẳng $AB$.
\choice{\True $x+y+2z-3=0$}{$x+y+2z-6=0$}{$x+3y+4z-7=0$}{$x+3y+4z-26=0$}
\loigiai{$\overrightarrow{AB}=(1;1;2)$ là một vectơ pháp tuyến của $(P)$. Phương trình $(P)$: $1(x-0)+1(y-1)+2(z-1)=0\Leftrightarrow x+y+2z-3=0$.}
\end{ex}
